\section{Conclusão}
\label{sec:conclusao}

Sob um surto de frente rápida, o enrolamento de um gerador síncrono não se comporta como uma indutância
concentrada: no instante inicial ele é uma rede capacitiva, e a tensão se distribui de forma
fortemente não uniforme, concentrada nas primeiras espiras. O grau de concentração é fixado por um
único número, \(\adist = \sqrt{\Cg/\Cs}\); para a máquina estudada, \(\adist = 5\) implica que a
entrada é solicitada por cerca de cinco vezes o valor de uma distribuição uniforme.

Como a barra do estator fica embutida numa ranhura de ferro aterrado, a capacitância para terra \Cg{} é
grande — o que torna \(\adist\) elevado e a concentração severa. Esse é o mecanismo físico que justifica
as práticas de proteção contra surto nos terminais de máquinas rotativas: para-raios e, sobretudo,
\emph{capacitores de surto}, que reduzem a inclinação da frente que chega ao enrolamento e, portanto, a
concentração inicial.

Para o projeto de isolamento, a conclusão prática é que \emph{o envelope de tensão durante o
transitório} — e não o regime permanente — dimensiona a solicitação: as primeiras espiras concentram a
tensão no instante do surto, enquanto nós internos podem exceder a tensão de entrada durante a
redistribuição (a simulação indica pico em torno de \SI{1108}{V} a 25\,\% do enrolamento, para uma
entrada de \SI{1}{kV}). A condição de contorno é decisiva: com o neutro solidamente aterrado, o terminal
final é preso a zero e a solicitação relevante é \emph{interna} — diferentemente do caso de neutro
isolado, em que a reflexão elevaria a tensão na extremidade.

As afirmações deste relatório repousam sobre duas verificações independentes — a solução analítica
\(\sinh\) para a distribuição inicial e a simulação cruzada em ATP/EMTP para a evolução completa — que
concordam com o modelo dentro de margens estreitas. Como limitações, o modelo usa \Nsec{} seções
uniformes e não representa a transição ranhura\,–\,cabeça de bobina, o acoplamento mútuo entre bobinas
nem perdas dependentes da frequência; ainda assim, captura o mecanismo essencial: a distribuição inicial
não uniforme governada por \(\adist\) e sua relaxação no tempo.
